% Источники: 02_Общие_соглашения/Глоссарий.md
% Генератор: docs/build_pdf.py
% Шаблон печатного представления. Содержимое подставляет docs/build_pdf.py.
\documentclass[12pt,a4paper]{article}
\usepackage{iftex}
\ifPDFTeX
  \usepackage[utf8]{inputenc}
  \usepackage[T2A]{fontenc}
  \usepackage{newtxtext,newtxmath}
\else
  \usepackage{fontspec}
  \IfFontExistsTF{Times New Roman}
    {\setmainfont{Times New Roman}}
    {\setmainfont{Libertinus Serif}}
  \IfFontExistsTF{Arial}{\setsansfont{Arial}}{\setsansfont{Latin Modern Sans}}
  \IfFontExistsTF{Menlo}{\setmonofont{Menlo}[Scale=MatchLowercase]}
    {\setmonofont{Latin Modern Mono}[Scale=MatchLowercase]}
\fi
\usepackage[russian]{babel}
\usepackage{geometry}
\geometry{left=25mm,right=15mm,top=20mm,bottom=20mm}
\usepackage{setspace}
\onehalfspacing
\usepackage{indentfirst}
\setlength{\parindent}{1.25cm}
\setlength{\parskip}{2pt}
\setlength{\emergencystretch}{3em}
\usepackage{graphicx,booktabs,longtable,array,tabularx,calc,etoolbox}
\usepackage{enumitem}
\setlist[itemize]{leftmargin=1.25cm,itemsep=2pt,topsep=4pt}
\setlist[enumerate]{leftmargin=1.25cm,itemsep=2pt,topsep=4pt}
\newcommand{\tightlist}{\setlength{\itemsep}{0pt}\setlength{\parskip}{0pt}}
\AtBeginEnvironment{longtable}{\small\singlespacing}
\setlength{\LTpre}{8pt}
\setlength{\LTpost}{8pt}
\setlength{\tabcolsep}{5pt}
\setlength{\LTleft}{0pt}
\setlength{\LTright}{0pt}
\usepackage{xurl}
\usepackage[unicode,colorlinks=true,linkcolor=black,urlcolor=blue]{hyperref}
\usepackage{bookmark}
\hypersetup{pdftitle={Глоссарий},pdfauthor={Анохин Г.А.},pdfsubject={Криптотренажёр mephi-12}}
\urlstyle{same}
\pagestyle{plain}
\setcounter{secnumdepth}{3}
\setcounter{tocdepth}{2}

\begin{document}
\hypersetup{pageanchor=false}
\begin{titlepage}
  \thispagestyle{empty}
  \centering
  {\large НИЯУ МИФИ\par}
  \vspace*{3.5cm}
  {\LARGE\bfseries Глоссарий\par}
  \vspace{0.5cm}
  {\large Глоссарий\par}
  \vspace{0.6cm}
  {\large Платформа для изучения и практики курса по криптографии\par}
  {\large Криптотренажёр mephi-12\par}
  \vspace{0.7cm}
  {\large Версия 0.3\par}
  {\large Редакция от 07.11.2026\par}
  \vspace{1cm}
  \begin{tabularx}{\textwidth}{@{}p{3.4cm}Xp{4cm}@{}}
    \toprule
    Роль & ФИО & Подпись / Дата \\
    \midrule
    Разработал & Анохин Г.А. & \\
    \addlinespace[6pt]
    Согласовано & & \\
    \addlinespace[6pt]
    Утверждено & & \\
    \bottomrule
  \end{tabularx}
  \vfill
  {\large Москва, 2026\par}
\end{titlepage}
\hypersetup{pageanchor=true}

\section*{Лист регистрации изменений}
\addcontentsline{toc}{section}{Лист регистрации изменений}
\begingroup\small\singlespacing
\noindent
\begin{tabularx}{\textwidth}{@{}p{1.2cm}p{2cm}p{2.5cm}Xp{1.8cm}@{}}
\toprule
Версия & Дата & Автор & Содержание изменения & Утвердил \\
\midrule
0.3 & 07.11.2026 & Анохин Г.А. & Уточнены правила и контракты; расширено описание сценариев; подготовлена печатная редакция. & \\
\bottomrule
\end{tabularx}
\endgroup
\clearpage
\begingroup\small\singlespacing
\tableofcontents
\endgroup
\clearpage

\hypertarget{глоссарий}{}

Версия: 0.3. Дата: 07.11.2026. Ответственный: Анохин Г.А. Определения относятся к рабочей редакции требований.

{\def\LTcaptype{none} % do not increment counter
\begin{longtable}[]{@{}
  >{\raggedright\arraybackslash}p{(\linewidth - 2\tabcolsep) * \real{0.3200}}
  >{\raggedright\arraybackslash}p{(\linewidth - 2\tabcolsep) * \real{0.6800}}@{}}
\toprule\noalign{}
\begin{minipage}[b]{\linewidth}\raggedright
Термин
\end{minipage} & \begin{minipage}[b]{\linewidth}\raggedright
Значение
\end{minipage} \\
\midrule\noalign{}
\endhead
\bottomrule\noalign{}
\endlastfoot
Криптотренажёр & Учебная система для выдачи криптографических задач и проверки ответов \\
Пользователь & Человек, работающий через свою учётную запись \\
Группа & Объединение пользователей с доступными ему конфигурациями задач \\
Конфигурация & Именованное описание состава задач; доступность определяется для группы \\
Сессия задач, набор задач & Набор задач одного пользователя с общим сроком и итоговым состоянием \\
Задача & Одно учебное задание с условием, типом и состоянием внутри сессии \\
Предметная проверка & Проверка правильности решения по правилам соответствующего криптографического домена \\
Попытка & Предметная проверка, завершённая фиксацией вердикта «верно» или «неверно». Ошибка формата и технический сбой без вердикта попытку не расходуют \\
Итоговое состояние & \texttt{SOLVED} или \texttt{FAILED}; повторная отправка ответа его не изменяет \\
\texttt{NEW} задачи & Итог задачи ещё не определён; возможность приёма ответа дополнительно зависит от срока и наличия уже принятой проверки \\
\texttt{SOLVED} задачи & Решение проверено и признано верным \\
\texttt{FAILED} задачи & Решение признано неверным либо задача завершена по просрочке \\
\texttt{NEW} сессии & Есть задачи без итогового состояния; после срока новые ответы уже не принимаются \\
\texttt{SOLVED} сессии & Все задачи получили \texttt{SOLVED} \\
\texttt{FAILED} сессии & Все задачи завершены, хотя бы одна получила \texttt{FAILED} \\
Момент получения ответа & Зафиксированное сервером время получения полного запроса с ответом; время на устройстве пользователя не используется \\
Своевременный ответ & Ответ, полученный сервером строго до общего срока сессии; предметная проверка может закончиться после срока \\
Ошибка формата & Запрос нельзя передать предметной проверке из-за синтаксиса, типа данных или отсутствующего обязательного поля \\
Неверный ответ & Корректно оформленный ответ, который предметная проверка признала неправильным \\
Технический сбой проверки & Проверка не завершилась вердиктом из-за отказа программного компонента или его зависимости \\
Контракт & Обязательства участников взаимодействия: входы, результаты, условия и отказы \\
SR & Системное требование к наблюдаемому поведению или свойству системы \\
SWR & Программное требование, обеспечивающее системное требование \\
Трассировка & Явная связь требований разных уровней и описаний их проверок \\
ФИ & Формальная инспекция документа с назначенными ролями, замечаниями и решением \\
\end{longtable}
}

Точные правила срока и попытки заданы в \href{../domains/\%D0\%A1\%D0\%B5\%D1\%81\%D1\%81\%D0\%B8\%D0\%B8_\%D0\%B7\%D0\%B0\%D0\%B4\%D0\%B0\%D1\%87.pdf\#nameddest=\%D0\%BA\%D0\%B0\%D1\%80\%D1\%82\%D0\%BE\%D1\%87\%D0\%BA\%D0\%B0-\%D0\%B4\%D0\%BE\%D0\%BC\%D0\%B5\%D0\%BD\%D0\%B0-\%D1\%81\%D0\%B5\%D1\%81\%D1\%81\%D0\%B8\%D0\%B8-\%D0\%B7\%D0\%B0\%D0\%B4\%D0\%B0\%D1\%87}{карточке сессий}. Обозначения требований определены в \href{\%D0\%9C\%D0\%B5\%D1\%82\%D0\%BE\%D0\%B4\%D0\%B8\%D0\%BA\%D0\%B0_\%D0\%B8\%D0\%B4\%D0\%B5\%D0\%BD\%D1\%82\%D0\%B8\%D1\%84\%D0\%B8\%D0\%BA\%D0\%B0\%D1\%82\%D0\%BE\%D1\%80\%D0\%BE\%D0\%B2.pdf\#nameddest=\%D0\%BC\%D0\%B5\%D1\%82\%D0\%BE\%D0\%B4\%D0\%B8\%D0\%BA\%D0\%B0-\%D0\%B8\%D0\%B4\%D0\%B5\%D0\%BD\%D1\%82\%D0\%B8\%D1\%84\%D0\%B8\%D0\%BA\%D0\%B0\%D1\%82\%D0\%BE\%D1\%80\%D0\%BE\%D0\%B2}{методике}.


\end{document}
